\chapter{Escrever com IA sem Perder a Tua Voz}
\label{ch:escrita}

\ficcao A biblioteca do bairro mudou a hora de uma oficina. O primeiro rascunho anunciou uma «experiência imperdível». O segundo garantiu lugar a toda a gente. A nota da organização dizia apenas que a oficina começava às dez e que as inscrições dependiam das vagas. Qual dos avisos deixarias na porta?

\section{Uma revisão que conserva os factos}

Começa por separar o que o leitor precisa de fazer daquilo que o texto precisa de dizer. No aviso da biblioteca, a pessoa quer saber se ainda pode inscrever-se e a que horas deve chegar. A voz pode ser informal ou institucional. A hora e a condição de acesso continuam a ser as mesmas.

\begin{quote}
\small
\textbf{Nota original do exercício.} Oficina no sábado, às 10h, na biblioteca. Inscrições até quinta-feira, sujeitas a vagas.

\textbf{Aviso possível.} A oficina começa às 10h de sábado, na biblioteca. Podes inscrever-te até quinta-feira, enquanto houver vagas.
\end{quote}

A comparação tem duas passagens. Primeiro confronta dia e hora, local, prazo e condição com a nota. Depois lê o aviso em voz alta e verifica se o pedido fica claro na primeira leitura. Trocar «enquanto houver vagas» por «garante o teu lugar» mudaria o compromisso, mesmo que a frase ficasse mais fluida.

\begin{promptlabbox}{Romper o «muro bege»}
\textbf{Objetivo e contexto:} Rever um parágrafo com os mesmos factos. Usa a nota da oficina acima ou um texto fictício com condições explícitas.

Escreve uma versão para um aviso afixado na porta e outra para uma mensagem a quem se quer inscrever. Indica o público e o espaço disponível. Podes fazer as duas versões à mão ou pedir rascunhos a um modelo.

\textbf{Resultado esperado e validação:} Sublinha dia e hora, local, prazo e condição na nota. Confere se são recuperáveis em ambas as versões e se cada leitor percebe o próximo passo. Se comparares a legibilidade, oculta qual das versões veio do modelo.

\textbf{Falhas possíveis:} Prometer lugar, omitir a condição de disponibilidade ou acrescentar um contacto inexistente.
\end{promptlabbox}

\section{Passar de frases vagas a ações}

«Vamos potenciar a comunicação» pede ao leitor que adivinhe o que muda. «À segunda-feira, publicamos as datas da semana no quadro da entrada» indica uma ação, um momento e um lugar. Se esses detalhes ainda estiverem por decidir, assinala-os no rascunho e pede confirmação. A revisão que inventa um responsável resolve a frase e cria um problema novo.

\begin{lstlisting}
Revê o texto para o público indicado.
Mantém datas, valores e condições das notas.
Assinala cada promessa nova como ponto a confirmar.
Por fim, indica o que o leitor deve conseguir fazer.
\end{lstlisting}

Podes dar um exemplo do tom que procuras. Descreve o que o torna adequado, como frases de extensão variada, tratamento na segunda pessoa ou explicação de siglas à primeira ocorrência. A escolha de ritmo depende da tarefa. Um aviso de alteração de horário pode terminar quando a hora estiver clara. Um relato pode precisar de uma pausa antes de chegar à resposta.

\section{Descrever uma voz com exemplos próprios}

Para adaptar o texto a públicos diferentes, usa amostras tuas ou materiais que possas utilizar. Pede uma descrição de traços concretos e testa-os num assunto novo. Na versão final, verifica se os exemplos deixaram marcas que pertenciam à situação anterior, como uma morada, um nome ou uma piada que só fazia sentido naquela conversa.

\begin{promptlabbox}{Descrever e testar a tua voz}
\textbf{Objetivo e contexto:} Extrair traços de dois textos teus que possas partilhar e aplicá-los a um terceiro assunto.

Pede uma ficha com tom, ritmo, vocabulário e estrutura. Escolhe três traços que reconheças e dá ao modelo uma nota factual diferente das amostras. Pede um rascunho e lê-o em voz alta. Se preferires trabalhar sem modelo, usa a ficha como guião para a tua própria revisão.

\textbf{Resultado esperado e validação:} O novo texto cumpre o objetivo e conserva os factos da nota. Localiza um exemplo de cada traço que escolheste e verifica se as amostras não lhe emprestaram factos ou frases inadequadas.

\textbf{Falhas possíveis:} Descrever a voz só com adjetivos genéricos, copiar dados de uma amostra ou interpretar semelhança como autorização para falar por outra pessoa.
\end{promptlabbox}

\section{Exemplos que deixam ver o mecanismo}

A expressão inglesa \emph{show, don't tell} recomenda mostrar uma consequência concreta. Para explicar pesquisa documental, podes acompanhar uma pessoa que encontra duas instruções contraditórias e precisa de verificar qual está em vigor. Identifica a cena como fictícia. Para descrever um caso real, localiza os documentos, a data e o resultado observado. Os dois textos podem ter a mesma estrutura narrativa e exigir tipos diferentes de verificação.

\begin{promptlabbox}{Testar se um exemplo ajuda a agir}
\textbf{Objetivo e contexto:} Comparar duas explicações do mesmo conceito, com factos e extensão semelhantes.

Escolhe um procedimento que conheças. Numa versão, explica-o sem exemplo. Na outra, acompanha uma pessoa numa situação fictícia em que tenha de aplicar a regra. Dá a um leitor uma situação diferente e pergunta que passo faria primeiro.

\textbf{Resultado esperado e validação:} A resposta do leitor indica se conseguiu aplicar a explicação ao caso novo. Regista também o que o confundiu. Uma avaliação informal informa esta revisão, com alcance limitado à tarefa e aos leitores envolvidos.

\textbf{Falhas possíveis:} Acrescentar números sem fonte, transformar o exemplo num testemunho real ou escolher só situações que favorecem uma das versões.
\end{promptlabbox}

\section{Estilo, proveniência e decisões editoriais}

A \emph{perplexidade} exprime quão previsível é um texto segundo um modelo avaliador e a sua tokenização. O termo \emph{burstiness} é usado informalmente para descrever variação de ritmo. Um texto convencional de uma pessoa pode ser previsível. Sadasivan e colegas observaram quedas na deteção para vários detetores quando aplicaram paráfrases aos textos do estudo~\citep{sadasivan2023}. Se estiver em causa a autoria, considera o histórico de trabalho, as fontes e as regras do contexto. Um valor calculado sobre o texto é uma evidência parcial.

Para aprovar uma peça, regista a fonte dos factos, a condição que precisa de sobreviver à revisão e quem pode assumir o compromisso. Na comunicação profissional, a mesma revisão estende-se a decisões, responsáveis e prazos. O capítulo~\ref{ch:negocios} mostra onde cada um deve aparecer.
